%%%%%%%%%%%%%%%%%%%%%%%%%%%%%%%%%%%%%%%%%%%%%%%%%%%%%%%%%%%%%%%%%%%%%%%%
%                                                                      %
%     File: Thesis_Resumo.tex                                          %
%     Tex Master: Thesis.tex                                           %
%                                                                      %
%     Author: Andre C. Marta                                           %
%     Last modified :  4 Mar 2024                                      %
%                                                                      %
%%%%%%%%%%%%%%%%%%%%%%%%%%%%%%%%%%%%%%%%%%%%%%%%%%%%%%%%%%%%%%%%%%%%%%%%

\section*{Resumo}

% Add entry in the table of contents as section
\addcontentsline{toc}{section}{Resumo}

A utilização de próteses em animais tem vindo a crescer, impulsionada pelos avanços na medicina veterinária e engenharia biomédica. No entanto, a avaliação da adaptação e eficácia destas próteses continua a basear-se predominantemente em observações qualitativas, dependentes da interpretação subjetiva de médicos veterinários, o que limita a consistência e reprodutibilidade dos resultados.

Este trabalho propõe o desenvolvimento de uma prótese instrumentada para monitorização do comportamento locomotor de cães, através da aquisição de dados biomecânicos relevantes. A solução integra sensores inerciais e de força, permitindo medir parâmetros como aceleração, orientação e padrões de carga durante a marcha, recorrendo a um sistema embebido de baixo consumo com armazenamento local de dados.

Dado que a instrumentação se encontra limitada ao membro protético, a avaliação da adaptação não é realizada através de métricas clássicas de simetria entre membros. Em alternativa, são propostas métricas proxy baseadas na consistência temporal do movimento, variabilidade dos sinais e padrões de distribuição de força, permitindo inferir o nível de adaptação do animal.

Com base nestes dados, será explorada a utilização de modelos de inteligência artificial para classificação automática do estado de adaptação à prótese, enquadrando o problema como uma tarefa de classificação supervisionada. A seleção dos modelos será realizada numa fase posterior, em função das características dos dados adquiridos.

Espera-se que esta abordagem contribua para a transição de métodos qualitativos para métodos quantitativos e automatizados na avaliação de próteses, melhorando o processo de adaptação e otimização destes dispositivos. Adicionalmente, os resultados poderão ter impacto na área das próteses humanas, ao introduzir metodologias de avaliação mais rigorosas e baseadas em dados.
\vfill

\textbf{\Large Palavras-chave:} próteses caninas, análise de marcha
, sistemas embebidos, sensores inerciais, inteligência artificial
